% =====================================================================
\section{Proofs of the SO(3) identities}
\label{app:so3}
% =====================================================================

\textbf{\eqref{eq:id1}.} $\skw{\bm a}^\top=-\skw{\bm a}$ by inspection of \eqref{eq:hat}. $\skw{\bm a}\bm a=\bm a\times\bm a=0$, and
$\skw{\bm a}\bm b=\bm a\times\bm b=-\bm b\times\bm a=-\skw{\bm b}\bm a$.

\textbf{\eqref{eq:id2}.} For any $\bm b$: $R\skw{\bm a}R^\top\bm b=R\big(\bm a\times R^\top\bm b\big)=(R\bm a)\times(RR^\top\bm b)=(R\bm a)\times\bm b$,
because a rotation preserves cross products ($R(\bm x\times\bm y)=R\bm x\times R\bm y$ for $\det R=+1$).

\textbf{\eqref{eq:id3}.} Split $M=M_s+M_a$ into symmetric and skew parts. The trace of (symmetric $\times$ skew) is zero, so
$\tr(M\skw{\bm\eta})=\tr(M_a\skw{\bm\eta})$. With $M_a=\skw{\bm m}$, $\bm m=\tfrac12(M-M^\top)^\vee$, a direct calculation gives
$\tr(\skw{\bm m}\skw{\bm\eta})=-2\,\bm m^\top\bm\eta$. Hence $\tr(M\skw{\bm\eta})=-\bm\eta^\top(M-M^\top)^\vee$.

\textbf{Trace and angle.} For $R=\Exp(\theta\bm n)$, \eqref{eq:exp} gives $\tr R=3-(1-\cos\theta)\cdot2=1+2\cos\theta$ (since
$\tr\skw{\bm n}=0$ and $\tr\skw{\bm n}^2=-2$). This is used for $\Psi=1-\cos\varphi$ and for $\Log$.

% =====================================================================
\section{Matrix-calculus rules used}
\label{app:calc}
% =====================================================================

For vectors $\vx,\bm b$ and matrices $M$ (symmetric where stated), $K$, $N$:
\[
\frac{\partial(\bm b^\top\vx)}{\partial\vx}=\bm b,\qquad
\frac{\partial(\vx^\top M\vx)}{\partial\vx}=2M\vx\ (M=M^\top),
\]
\[
\frac{\partial\,\tr(KN)}{\partial K}=N^\top,\qquad
\frac{\partial\,\tr(KMK^\top)}{\partial K}=2KM\ (M=M^\top).
\]
The first two give the LQR gain (Model~2, block~4) and the QP gradient $H\mathbf U+\bm f$ (Model~2, block~8). The last two give the
Kalman gain (Model~1, block~4).

\textbf{Least squares.} $\min_G\|Y-G\Omega\|_F^2$ has gradient $-2(Y-G\Omega)\Omega^\top$. It is zero at
$G=Y\Omega^\top(\Omega\Omega^\top)^{-1}=Y\Omega^\dagger$ when $\Omega$ has full row rank (DMDc, Model~2 block~2).

% =====================================================================
\section{References}
% =====================================================================

\begin{enumerate}[label={[\arabic*]}, leftmargin=*]
\item T.~Lee, M.~Leok, N.\,H.~McClamroch, ``Geometric tracking control of a quadrotor UAV on SE(3)'',
      \emph{IEEE CDC}, 2010 (arXiv:1003.2005). Basis of Model 1, blocks 5--7.
\item J.~Sol\`a, ``Quaternion kinematics for the error-state Kalman filter'', arXiv:1711.02508, 2017. Basis of Model 1, block 4.
\item J.\,L.~Proctor, S.\,L.~Brunton, J.\,N.~Kutz, ``Dynamic mode decomposition with control'',
      \emph{SIAM J.\ Applied Dynamical Systems} 15(1), 2016. Model 2, block 2.
\item D.\,P.~Bertsekas, \emph{Dynamic Programming and Optimal Control}, Vol.~I, Athena Scientific. Model 2, blocks 3--6.
\item J.\,B.~Rawlings, D.\,Q.~Mayne, M.\,M.~Diehl, \emph{Model Predictive Control: Theory, Computation, and Design}, 2nd ed., 2017.
      Model 2, blocks 7--9.
\item T.~Qin, P.~Li, S.~Shen, ``VINS-Mono: a robust and versatile monocular visual-inertial state estimator'',
      \emph{IEEE T-RO}, 2018.
\item ArduPilot SITL and Gazebo Iris model documentation; project parameters from
      \code{controller.yaml} (package \code{drones\_controller}, folder \code{config}).
\end{enumerate}

% =====================================================================
\section{Glossary}
\label{app:glossary}
% =====================================================================

\begin{center}\small
\begin{tabularx}{\linewidth}{@{}l X@{}}
\toprule
\textbf{Term} & \textbf{Meaning}\\
\midrule
ArduPilot SITL & ArduPilot flight software running on a PC (Software-In-The-Loop), flying a simulated drone\\
Attitude & orientation of the drone; here the rotation matrix $R$\\
Body frame & axes fixed to the drone: $x$ forward, $y$ left, $z$ up (FLU)\\
Covariance & matrix of variances/co-variations of errors; the EKF's ``how unsure am I''\\
DARE & Discrete Algebraic Riccati Equation; gives the infinite-horizon LQR solution\\
DMDc & Dynamic Mode Decomposition with control; fits $\vx_{k+1}=A\vx_k+B\vu_k$ to data by least squares\\
EKF & Extended Kalman Filter: Kalman filter applied to a nonlinear model through linearization\\
ENU & East--North--Up world frame used by ArduPilot's ROS 2 topics\\
Error-state & the filter tracks the small error $\delta\vx$, not the full state, so attitude is handled correctly\\
Gazebo & 3D physics simulator in which the Iris quadrotor flies\\
Gimbal lock & singularity of Euler angles at $\pm90^\circ$ pitch; avoided by working with $R$ directly\\
IMU & Inertial Measurement Unit: gyroscope (rates) and accelerometer\\
Innovation & measurement minus what the filter expected to measure\\
Jacobian & matrix of partial derivatives; the ``slope'' of a vector function\\
Kalman gain $K$ & how strongly the estimate is pulled towards a measurement\\
LQR & Linear Quadratic Regulator: optimal linear feedback $\vu=-K\vx$\\
Lyapunov function & energy-like quantity that never increases; proves stability\\
MPC & Model Predictive Control: repeatedly solve a finite-horizon optimal problem with constraints\\
QP & Quadratic Program: minimise a quadratic cost subject to linear constraints\\
Receding horizon & apply only the first planned input, then re-plan one step later\\
RK4 & 4th-order Runge--Kutta: accurate way to step a differential equation forward\\
ROS 2 & Robot Operating System; our controller node talks to ArduPilot through it\\
SO(3) & the set of all 3D rotation matrices ($R^\top R=I$, $\det R=1$)\\
VIO & Visual-Inertial Odometry: estimating motion by fusing camera and IMU\\
ZOH & Zero-Order Hold: input kept constant between samples\\
\bottomrule
\end{tabularx}
\end{center}
